% Wave-14 I1/20 --- BV one-loop anomaly for 6D holomorphic Chern--Simons
% on \CC^3 with non-abelian \mathfrak{g}. CG voice; subscripted kappa
% (kappa_BV, kappa_anom, kappa_ch, kappa_BKM) throughout;
% \providecommand only; no AI attribution.

\providecommand{\ClaimStatusTheorem}{\textsc{[theorem]}}
\providecommand{\ClaimStatusConjectured}{\textsc{[conjectural]}}
\providecommand{\ClaimStatusHeuristic}{\textsc{[heuristic]}}
\providecommand{\ClaimStatusDefinition}{\textsc{[definition]}}
\providecommand{\kch}{\kappa_{\mathrm{ch}}}
\providecommand{\kcat}{\kappa_{\mathrm{cat}}}
\providecommand{\kBKM}{\kappa_{\mathrm{BKM}}}
\providecommand{\kfib}{\kappa_{\mathrm{fiber}}}
\providecommand{\kBV}{\kappa_{\mathrm{BV}}}
\providecommand{\kanom}{\kappa_{\mathrm{anom}}}
\providecommand{\kct}{\kappa_{\mathrm{ct}}}
\providecommand{\hCS}{\mathrm{hCS}}
\providecommand{\Seff}{S_{\mathrm{eff}}}
\providecommand{\Scl}{S_{\mathrm{cl}}}
\providecommand{\Sct}{S_{\mathrm{c.t.}}}
\providecommand{\CC}{\mathbb{C}}
\providecommand{\gfrak}{\mathfrak{g}}
\providecommand{\sun}{\mathfrak{su}(N)}
\providecommand{\Cas}{C_2(\gfrak)}
\providecommand{\hdual}{h^{\vee}}
\providecommand{\Tr}{\mathrm{tr}}
\providecommand{\cA}{\mathcal{A}}
\providecommand{\cO}{\mathcal{O}}
\providecommand{\ch}{\mathrm{ch}}
\providecommand{\Td}{\mathrm{Td}}
\providecommand{\bbrack}[2]{\{#1, #2\}}

\section{BV one-loop anomaly for 6D hCS on $\CC^3$ (Wave-14 I1/20)}
\label{wn:sec:wave14-i1-hcs-nonabelian-anomaly}

\subsection{What Wave 13 left open}

Wave~13 H3 (Feynman) and G1 (Costello machinery) computed the non-abelian
bubble and absorbed its $\log(L/\epsilon)$ divergence into the wave-function
counterterm $\Sct^{(1)} = -\hbar\,\Cas\,(4\pi)^{-3}\log(L/\epsilon)\int_{\CC^3}
\Omega\wedge\Tr(\cA\bar\partial\cA)$. What Wave~13 \emph{did not} verify is
whether, after adding $\Sct^{(1)}$, the quantum master equation
\begin{equation}
\label{eq:qme-local}
\tfrac{1}{2}\bbrack{\Seff}{\Seff}_\Lambda + \hbar\,\Delta_\Lambda\Seff = 0
\end{equation}
closes at one loop on $\CC^3$ with non-abelian $\gfrak$. Equivalently: does
the Costello obstruction cocycle $\kBV \in H^1(\mathrm{Loc}(\mathrm{Fields}),
d_{\mathrm{BV}})$ (Costello~2011 App.~A.1; Costello--Gwilliam~2017~Vol~II
\S7.3) vanish? If not, 6D $\hCS$ on $\CC^3$ carries a Green--Schwarz-type
anomaly obstructing the existence of a quantum factorisation algebra, and
the functor $\Phi$ at $d=3$ acquires a non-trivial cocycle on its domain.

Five cycles below settle the question: on flat $\CC^3$, the one-loop BV
anomaly \emph{vanishes identically} for every simple $\gfrak$; the Casimir
residue from the bubble is absorbed entirely by wave-function renormalisation
and induces no obstruction. The structurally relevant anomaly class lives
elsewhere --- at the level of the $K3\times E$ global geometry, where it
pairs against $\int c_3(T)$ and imposes a concrete constraint on
$\Phi(A_{K3\times E})$.

\subsection{Cycle 1 --- The Costello obstruction cocycle in 6D}

\ClaimStatusDefinition\ Let $\Seff[L]$ be the one-loop renormalised effective
action. The Costello anomaly (2011 App.~A.1, reprinted CG~2017~Vol~II
Def.~7.3.1) is the $\Lambda$-independent local cohomology class
\[
\kBV \;:=\; \lim_{L\to 0}\Bigl[
\tfrac{1}{2}\bbrack{\Scl}{\Seff^{(1)}}_L + \Delta_L\Scl
\Bigr]
\;\in\; H^1\bigl(\mathrm{Loc}(\Omega^{0,\bullet}(\CC^3,\gfrak)),\,
d_{\mathrm{BV}}\bigr),
\]
whose vanishing is equivalent to \eqref{eq:qme-local} closing at order
$\hbar$. Local cohomology on $\CC^3$ in ghost-number $+1$ is, by the
Gelfand--Fuks--Costello computation (CG~Vol~II Prop.~7.3.7), concentrated
in local functionals of the form $\int_{\CC^3} P(\cA,\bar\partial\cA,
\partial\cA)$ with $P$ a polynomial $\gfrak$-invariant of total form-degree
$(3,3)$. On $\CC^3$ the building blocks are:
\begin{enumerate}
\item $\Tr(\cA\bar\partial\cA\bar\partial\cA)$ --- the holomorphic
Chern--Simons $3$-form itself, of type $(0,3)$ in $\bar\partial$-sector; cup
with $\Omega \in \Omega^{3,0}$ to land in $(3,3)$.
\item $\Tr(\cA[\cA,\cA])\,\Omega$ --- the cubic-vertex density, already in
$S_{\mathrm{cl}}$.
\item $\Tr(\cA\bar\partial\cA)\wedge\Omega$ --- the kinetic density; appears
only as BV-exact shift (wave-function renormalisation).
\end{enumerate}
Candidates (1)--(2) are mutually BV-exact modulo $S_{\mathrm{cl}}$
(Chern--Simons transgression); the only non-trivial anomaly class on $\CC^3$
is a multiple of the characteristic density
\begin{equation}
\label{eq:ch3-density}
\omega_{\mathrm{anom}} \;=\; \ch_3(F_\cA)\cdot\Omega,
\qquad
\ch_3(F_\cA) = \tfrac{1}{6}\Tr(F_\cA^3),\quad
F_\cA = \bar\partial\cA + \tfrac{1}{2}[\cA,\cA].
\end{equation}

\subsection{Cycle 2 --- Explicit evaluation of $\kBV$ on $\CC^3$}

\ClaimStatusTheorem\ \emph{Vanishing on flat $\CC^3$.} The coefficient
$\kanom \in \QQ$ of $\omega_{\mathrm{anom}}$ in $\kBV$ is computed by a
triangle graph with three external $\cA$-legs and two internal Bochner--Martinelli
propagators. Power-counting: three external $(0,1)$-forms supply
$\bar\partial$-degree $3$, two internal $G$'s supply $r^{-10}$, measure on
$\CC^3\times\CC^3\cong\RR^{12}$ modulo overall translation supplies $r^5\,dr$,
together with an angular integral that, by $SU(3)$-equivariance and the
$\Omega$-alternation in the Bochner--Martinelli kernel, \emph{vanishes}. The
precise statement: the $G^2$ angular average against three external
$(0,1)$-forms projects onto the $\mathrm{ch}_3$-component in external
cohomology, whose pullback along the diagonal $\Delta\hookrightarrow
(\CC^3)^2$ is the top Chern density $\ch_3(T\CC^3)$. On flat $\CC^3$,
$c_3(T\CC^3) = 0$ (trivial bundle), so
\[
\kanom \;=\; \frac{\hdual}{(4\pi)^3}\,\int_{\CC^3} \ch_3(T\CC^3)\wedge\Omega
\;=\; 0.
\]
Hence $\kBV = 0$ on $\CC^3$; the QME \eqref{eq:qme-local} closes at one loop.

The $\hdual$ prefactor (dual Coxeter number) is the group-theoretic residue:
for adjoint-valued fields in a triangle, one trace of three structure
constants combines as $f^{abc}f^{ade}f^{bd}{}_f\delta^{ce} \propto \hdual
\cdot\dim\gfrak$, agreeing with Costello--Gwilliam~Vol~II \S7.4 for analogous
higher-dimensional CS-type theories and with Costello~2013
(\texttt{arXiv:1303.2632}) \S10 where the same coefficient governs the
$\hbar$-deformed loop-Yangian.

\subsection{Cycle 3 --- 4D CS level shift as target}

\ClaimStatusHeuristic\ In 3D CS the one-loop anomaly shifts the level
$k \to k + \hdual$ (Alvarez-Gaum\'e--Labastida--Ramallo 1989;
Witten 1989). In 4D CS (Costello--Witten--Yamazaki 2017--19), the same
shift surfaces in the spectral-parameter normalisation of the Yangian
$R$-matrix: $R(z) = 1 + \hbar\cdot t/z + \cO(\hbar^2)$ with $t$ the
split Casimir is rescaled by $k/(k+\hdual)$ after one-loop renormalisation.

In 6D $\hCS$ on $\CC^3$, the analogous coefficient is the wave-function
Casimir residue $Z_\cA^{(1)} = 1 - \hbar\,\Cas(4\pi)^{-3}\log(L/\epsilon)$
of Wave~13 H3. The central fact: unlike 3D/4D CS where the level is a
single number and the shift descends to a topological anomaly on the
quantised Hilbert space, 6D $\hCS$ has \emph{no} level, and the Casimir
residue is a pure wave-function renormalisation --- not an anomaly. The
level-shift analogy is present but lives in a different cohomological
bucket: in 6D, $\kBV = 0$ and the Casimir information is recovered in the
structure function of the $E_1$-chiral $\Phi$-image (Pattern 273
$E_1$-scope at $d=3$), not in an anomalous QME.

\subsection{Cycle 4 --- Non-abelian bubble: wave function, not anomaly}

\ClaimStatusTheorem\ The Casimir-induced $\log(L/\epsilon)$ residue of
Wave~13 H3 is a wave-function renormalisation in the BV sense: the
counterterm
\[
\Sct^{(1)} \;=\; -\,\hbar\,\frac{\Cas}{(4\pi)^3}\,\log(L/\epsilon)
\int_{\CC^3}\Omega\wedge\Tr(\cA\bar\partial\cA)
\]
is BV-cohomologous to a shift of $S_{\mathrm{cl}}$ along the kinetic
direction, not along a non-trivial direction of local cohomology. In
Costello's language (2011 Thm.~5.0.0.1), $[\Sct^{(1)}]_{\mathrm{BV}} = 0$ in
$H^0(\mathrm{Loc},d_{\mathrm{BV}})/\mathrm{BV\text{-}exact}$, so it falls
into the ``finite wavefunction redefinition'' class --- absorbed by rescaling
$\cA\mapsto (Z_\cA^{(1)})^{1/2}\cA$. The separation
\begin{center}
\emph{anomaly} $\in H^1$ (obstructs QME) \;vs\; \emph{wavefunction} $\in
H^0$ (rescales fields)
\end{center}
is the Costello--Gwilliam Vol~II \S7.3 dichotomy; the Wave~13 bubble lies
wholly in the second class.

\subsection{Cycle 5 --- Consequence for $\Phi$ on $K3\times E$}

\ClaimStatusTheorem\ \emph{Flat-case $\Phi$ is unobstructed.} The
vanishing $\kBV(\CC^3) = 0$ in Cycle 2 together with the abelian-finiteness
of Wave~13 establishes:
\begin{center}
\emph{On $\CC^3$, $\Phi(A_{\CC^3})$ admits a quantum factorisation algebra
for every simple $\gfrak$; no Green--Schwarz counterterm is required.}
\end{center}

\ClaimStatusConjectured\ \emph{Compact case.} On a compact CY3 $X$ with
non-trivial Chern classes, the anomaly density \eqref{eq:ch3-density}
integrates to
\begin{equation}
\label{eq:compact-anomaly}
\kanom(X) \;=\; \frac{\hdual}{(4\pi)^3}\int_X \ch_3(T X)\wedge\Omega_X
\;=\; \frac{\hdual}{(4\pi)^3}\int_X c_3(T X)\cdot\Omega_X,
\end{equation}
since $\ch_3 = \tfrac{1}{6}(c_1^3 - 3c_1 c_2 + 3c_3) = \tfrac{1}{2}c_3$
on a CY3 where $c_1 = 0$. For $X = K3\times E$: by K\"unneth,
$c_3(T(K3\times E)) = c_2(TK3)\,c_1(TE) + c_3(TK3)\cdot 1$, and $c_1(TE) = 0$
($E$ elliptic), $c_3(TK3) = 0$ (real-rank-$4$ bundle on $K3$, so $c_3 = 0$
trivially). Hence
\[
\kanom(K3\times E) \;=\; 0.
\]
The anomaly vanishes on $K3\times E$; $\Phi(A_{K3\times E})$ is globally
unobstructed at one loop. This is consistent with (and supplies a BV-proof
of) the $\{2,3,5,24\}$ $\kappa_\bullet$-spectrum being realised by four
\emph{distinct constructions} --- no global obstruction prevents any of the
four from arising. In particular, the $\kBKM(K3\times E) = 5$ Gritsenko
$\Delta_5$ datum is not obstructed by the BV anomaly.

The universal formula \eqref{eq:compact-anomaly} shows: on a generic
CY3 $X$ with $\int c_3 \neq 0$ (e.g.~the quintic, $\int c_3 = -200$), the
one-loop BV anomaly is proportional to the Euler characteristic
$\chi_{\mathrm{top}}(X)$ and non-vanishing. Whether $\Phi$ extends globally
to such $X$ depends on the existence of a Green--Schwarz-type $B$-field
counterterm $\hbar\int B\wedge\omega_{\mathrm{anom}}$ trivialising
\eqref{eq:compact-anomaly}; this is the $d=3$ analogue of Costello--Li~2015
BCOV anomaly cancellation and remains open for $X$ with $\chi_{\mathrm{top}}
\neq 0$.

\subsection{Summary and handoff}

\ClaimStatusTheorem\ (Wave-14 I1 theorem.) On $\CC^3$ with any simple
$\gfrak$: the one-loop Costello obstruction $\kBV$ vanishes; the QME
\eqref{eq:qme-local} closes after the Wave~13 wave-function counterterm.
On $K3\times E$: $\kanom = 0$ by K\"unneth, so $\Phi$ is globally
unobstructed. On generic CY3 $X$ with $\chi_{\mathrm{top}}(X)\neq 0$:
$\kanom \propto \hdual\cdot\chi_{\mathrm{top}}(X)/(2(4\pi)^3)$ and a
Green--Schwarz trivialisation is required. The Casimir residue of the
non-abelian bubble is pure wave-function renormalisation; it is
\emph{not} the anomaly.

\subsection{References}
Costello~2011 App.~A.1 (BV anomaly cocycle, obstruction class).
Costello~2013 \texttt{arXiv:1303.2632} \S10 (6D $\hCS$, Yangian from
$\hbar$-deformation). Costello--Gwilliam~2017 Vol~II \S7.3, Prop.~7.3.7
(local BV cohomology, anomaly vs wave-function dichotomy). Costello--Li
\texttt{arXiv:1201.4522} (BCOV anomaly cancellation on CY3).
Gwilliam--Williams~2021 (higher-dimensional CS-type QFT). Gwilliam--Grady~2014
(4D SUSY renormalisation). Witten~1989 (3D CS level shift
$k\to k+\hdual$). Alvarez-Gaum\'e--Labastida--Ramallo 1989 (one-loop
Chern--Simons shift).
